\documentclass{article}
\usepackage{amsmath,amssymb}
\begin{document}
\section*{\oplus COUNTING}

\begin{tabular}{l c c c}
House & R & G & B \\
[[DIAGRAM: human figure]] & III & III & IIII \\
[[DIAGRAM: simplified human figure]] & I & I & NONE \\
\end{tabular}

The ``recipe'' to go up in value: $1 \to \mathrm{II} \to \mathrm{III} \to \mathrm{IIII} \to \mathrm{V} \to \mathrm{X} \to \mathrm{L}$ etc.

\section*{\oplus NUMBERS}

This can become hard to read.

\begin{tabular}{c c c c c}
I & II & III & & \\
IV & V & VI & VII & VIII \\
IX & X & XI & XII & XIII \\
\end{tabular}

The Romans ``solved'' this. Most hands have \texttt{[[DIAGRAM: human figure]]} |||| or \texttt{V} fingers.

\texttt{II} hands is \texttt{V} and \texttt{II} is \texttt{X}.

The \texttt{IV} \& \texttt{IX} idea is clever.

\texttt{||||} < \texttt{||||} in fingers \texttt{[[DIAGRAM: hand with 4 fingers]]} < \texttt{[[DIAGRAM: hand with 3 fingers]]}.

\texttt{III} is \texttt{V} less 1 finger.

\texttt{IX} is 1 less than \texttt{X}.

\end{document}
